\documentclass[a4paper,11pt]{report}

\usepackage[german]{babel}
\usepackage{csquotes}
\usepackage{amsmath,amssymb}
\usepackage[pdfusetitle]{hyperref}
\usepackage{subfiles}
\usepackage{xparse}
\usepackage{mathtools} % für \coloneqq
\usepackage{booktabs}  % für \toprule, \midrule, \bottomrule
\usepackage{tabularx}  % für tabularx-Umgebung
\usepackage{makecell}

% erst amsthm (falls formula-envs.tex es nicht selbst lädt; doppelt ist meist ok, aber sauber ist einmal)
\usepackage{amsthm}
\usepackage{xparse}
\usepackage{multirow}
\usepackage{microtype}
\hypersetup{
  hidelinks,
  pdftitle={Die Grundlagen der Mathematik - Gesamtband},
  pdfauthor={Martin Kunze}
}


\input{tex/impl/thmlookup.tex}
% TeX-Schnittstelle fuer den zentralen Bandgraphen in band-dependencies.tsv.
% xr-hyper liest die AUX-Dateien; dieselbe Vorgängerliste steuert die
% thmlookup-Registries und die externen PDF-Ziele.
\RequirePackage{xr-hyper}

\directlua{
  local path = kpse.find_file("tex/impl/band-dependencies.lua")
  if not path then
    error("cannot find tex/impl/band-dependencies.lua")
  end
  banddeps = dofile(path)
}

\newcommand{\BandSetupStandalone}[1]{%
  \directlua{banddeps.emit_external_documents([==[\luaescapestring{\detokenize{#1}}]==])}%
  \directlua{banddeps.setup_standalone([==[\luaescapestring{\detokenize{#1}}]==])}%
}

\newcommand{\BandStartInMain}[1]{%
  % Keep the result numbers identical to the standalone volume. Sequential
  % hyperlink names remain unique despite the chapter and title-page resets.
  \setcounter{chapter}{0}%
  \setcounter{section}{0}%
  \hypersetup{hypertexnames=false}%
  \directlua{banddeps.start_band_in_main([==[\luaescapestring{\detokenize{#1}}]==])}%
}

\input{tex/impl/formula-envs.tex}
\input{tex/impl/proof-tables.tex}
\input{tex/impl/commands.tex}

% Verhindert Seitenumbruch innerhalb eines Regelschemas
\newenvironment{RuleSchema}{\begin{samepage}}{\end{samepage}}

\newcounter{file}

% --- Warnschwellen ---
\hbadness=10000
\vbadness=10000
\hfuzz=1pt
\vfuzz=1pt


\begin{document}

\hypersetup{hypertexnames=false}
% longtable creates its own destinations independently of hypertexnames.
\renewcommand{\theHtable}{\FormulaBandID.\arabic{chapter}.\arabic{table}}
\renewcommand{\theHfigure}{\FormulaBandID.\arabic{chapter}.\arabic{figure}}
\title{Die Grundlagen der Mathematik\\Gesamtband}
\author{Martin Kunze}
\date{}
\maketitle
\setcounter{tocdepth}{3}
\tableofcontents
% The subfiles keep their own front matter when built independently.
% In the complete manuscript the shared title and contents appear once.
\let\maketitle\relax
\let\tableofcontents\relax
% Some subfiles repeat their title metadata in the document body. After
% report's \maketitle these commands are \relax; consume their arguments
% explicitly so the metadata cannot become an extra text page in main.
\def\title#1{}
\def\author#1{}
\def\date#1{}
\clearpage

\BandStartInMain{B01}
\subfile{Bd. 01 - Grundlagen der Logik}

\clearpage
\BandStartInMain{B02}
\subfile{Bd. 02 - Theoreme der Logik}

\clearpage
\BandStartInMain{B03}
\subfile{Bd. 03 - Mengenlehre}

\clearpage
\BandStartInMain{B04}
\subfile{Bd. 04 - Totale Relationen}

\clearpage
\BandStartInMain{B05}
\subfile{Bd. 05 - Funktionen}

\clearpage
\BandStartInMain{B06}
\subfile{Bd. 06 - Injektive Funktionen}

\clearpage
\BandStartInMain{B07}
\subfile{Bd. 07 - Surjektive Funktionen}

\clearpage
\BandStartInMain{B08}
\subfile{Bd. 08 - Bijektive Funktionen}

\clearpage
\BandStartInMain{B09}
\subfile{Bd. 09 - Auswahlprinzip}

\clearpage
\BandStartInMain{B10}
\subfile{Bd. 10 - Natürliche Zahlen}

\clearpage
\BandStartInMain{B11}
\subfile{Bd. 11 - Äquivalenzrelationen}

\clearpage
\BandStartInMain{B12}
\subfile{Bd. 12 - Halbordnungen}

\clearpage
\BandStartInMain{B13}
\subfile{Bd. 13 - Schranken, Infima und Suprema}

\clearpage
\BandStartInMain{B14}
\subfile{Bd. 14 - Paarinfima und Paarsuprema}

\clearpage
\BandStartInMain{B15}
\subfile{Bd. 15 - Totale Ordnungen}

\clearpage
\BandStartInMain{B16}
\subfile{Bd. 16 - Wohlordnungen und Auswahlaxiom}

\clearpage
\BandStartInMain{B17}
\subfile{Bd. 17 - Ganze Zahlen}

\clearpage
\BandStartInMain{B18}
\subfile{Bd. 18 - Rationale Zahlen}

\clearpage
\BandStartInMain{B19}
\subfile{Bd. 19 - Reelle Zahlen}

\clearpage
\BandStartInMain{B20}
\subfile{Bd. 20 - Endliche Mengen}

\clearpage
\BandStartInMain{B21}
\subfile{Bd. 21 - Folgen}

\clearpage
\BandStartInMain{B22}
\subfile{Bd. 22 - Gerichtete Graphen}

\clearpage
\BandStartInMain{B23}
\subfile{Bd. 23 - Ungerichtete Graphen}

\clearpage
\BandStartInMain{B24}
\subfile{Bd. 24 - Schlichte ungerichtete Graphen}

\clearpage
\BandStartInMain{B25}
\subfile{Bd. 25 - Zusammenhängende Graphen}

\clearpage
\BandStartInMain{B26}
\subfile{Bd. 26 - Bäume}

\clearpage
\BandStartInMain{B27}
\subfile{Bd. 27 - Endliche Wörter und Klammerungsbäume}

\clearpage
\BandStartInMain{B28}
\subfile{Bd. 28 - Halbgruppen}

\clearpage
\BandStartInMain{B29}
\subfile{Bd. 29 - Kommutative Halbgruppen}

\clearpage
\BandStartInMain{B30}
\subfile{Bd. 30 - Idempotente Halbgruppen}

\clearpage
\BandStartInMain{B31}
\subfile{Bd. 31 - Rechtecksbänder}

\clearpage
\BandStartInMain{B32}
\subfile{Bd. 32 - Kommutative idempotente Halbgruppen}

\clearpage
\BandStartInMain{B33}
\subfile{Bd. 33 - Halbgruppen mit Nullelement}

\clearpage
\BandStartInMain{B34}
\subfile{Bd. 34 - Linkskürzbare Halbgruppen}

\clearpage
\BandStartInMain{B35}
\subfile{Bd. 35 - Rechtskürzbare Halbgruppen}

\clearpage
\BandStartInMain{B36}
\subfile{Bd. 36 - Kürzbare Halbgruppen}

\clearpage
\BandStartInMain{B37}
\subfile{Bd. 37 - Endliche Halbgruppen}

\clearpage
\BandStartInMain{B38}
\subfile{Bd. 38 - Monoide}

\clearpage
\BandStartInMain{B39}
\subfile{Bd. 39 - Halbringe}

\clearpage
\BandStartInMain{B40}
\subfile{Bd. 40 - Gruppen und Ringe}

\clearpage
\BandStartInMain{B41}
\subfile{Bd. 41 - Halbverbände und Verbände}

\clearpage
\BandStartInMain{B42}
\subfile{Bd. 42 - Frankls Vermutung}

\clearpage
\BandStartInMain{B43}
\subfile{Bd. 43 - Metrische Räume und Vollständigkeit}

\end{document}
